% Generated draft: requires mathematical and source review.
\documentclass[11pt]{article}
\usepackage[margin=25mm]{geometry}
\usepackage{fontspec}
\setmainfont{Latin Modern Roman}
\usepackage{amsmath,amssymb,mathtools,unicode-math}
\setmathfont{Latin Modern Math}
\usepackage{enumitem,needspace,xurl,hyperref}
\setlength{\parindent}{0pt}
\setlength{\parskip}{5pt}
\setlength{\emergencystretch}{3em}
\newcommand{\N}{\mathbb{N}}
\newcommand{\Z}{\mathbb{Z}}
\newcommand{\Q}{\mathbb{Q}}
\newcommand{\R}{\mathbb{R}}
\newcommand{\C}{\mathbb{C}}
\newcommand{\F}{\mathbb{F}}
\newcommand{\A}{\mathbb{A}}
\newcommand{\PP}{\mathbb{P}}
\newcommand{\E}{\mathbb{E}}
\newcommand{\eps}{\varepsilon}
\newcommand{\dd}{\,\mathrm{d}}
\begin{document}
\section*{UnsolvedMath ID: 30006651}
\textbf{Source problem code:} catalogue entry
\subsection{Definitions and mathematical statement}
Let \(G\) be a finite or profinite group. For a prime number \(p\), a finite \(p\)-group has order a power of \(p\), and a pro-\(p\) group is an inverse limit of finite \(p\)-groups. The source considers a finite \(p\)-subgroup or pro-\(p\) subgroup \(H\leq G\).

Write
\[
\mathcal K=\operatorname{DPerm}(G;k)^\omega,
\qquad
X=\operatorname{Spec}(\mathcal K),
\]
following the source notation for the compact derived permutation modules and their tensor-triangular spectrum. Under the usual tensor-triangular interpretation, points of \(X\) are prime thick tensor ideals of \(\mathcal K\): a proper thick tensor ideal \(\mathcal P\) is prime when \(a\otimes b\in\mathcal P\) implies \(a\in\mathcal P\) or \(b\in\mathcal P\). Let
\[
X_{\mathrm{cl}}=\{x\in X:\overline{\{x\}}=\{x\}\}.
\]
The source does not specify the coefficient object \(k\), the precise construction of \(\operatorname{DPerm}(G;k)\), or the definition of periodicity. Accordingly, write \(\operatorname{Per}(X)\) for its intended periodic locus, with that definition still requiring recovery from the report.

\textbf{Conjecture 7, as supplied.} For every finite or profinite group \(G\), the intended assertion is
\[
\operatorname{Per}\!\left(\operatorname{Spec}\bigl(\operatorname{DPerm}(G;k)^\omega\bigr)\right)
=
\operatorname{Spec}\bigl(\operatorname{DPerm}(G;k)^\omega\bigr)\setminus X_{\mathrm{cl}}.
\]
The source also states the pointwise formulation, for the indicated subgroups \(H\leq G\),
\[
P(H,p)\text{ is periodic}\quad\Longleftrightarrow\quad p\text{ is periodic}.
\]
Here \(P(H,p)\) and its parameter are not defined in the supplied passage. In particular, the occurrence of \(p\) as a prime number in the subgroup hypothesis must be distinguished from whatever object is meant by ``\(p\) is periodic.'' No identification of these roles, or replacement by an unspecified prime ideal, is assumed. The range of the point parameter and the correspondence needed to justify the asserted equivalence with the global formulation remain unresolved. These displays preserve the requested target but do not constitute a fully specified standalone conjecture without the missing definitions and hypotheses.

\textbf{Supplied qualifications and reported special cases.} Remark 8 proposes following Example 2 and Remark 3, whose contents are not supplied. It reports that the line bundle denoted \(\Sigma 1\) is not ample, with the stated exception \(p\nmid\#G\), and that establishing ampleness of a suitable family is difficult. The meaning of this order condition for infinite profinite groups is not specified. The remark reports success for elementary abelian \(p\)-groups in [BG25], with associated global sections called permutation twisted cohomology. It attributes a generalization to \(p\)-groups to Miller [Mil25], without an ampleness result, but with injectivity of the comparison map to the Dirac scheme and control of sections after applying modular fixed points, described as \(p\)-local control. Theorem 9, attributed to [Gal25], reports that Conjecture 7 holds when \(G\) is a finite \(p\)-group or pro-\(p\) group; the passage also refers to [Gal25] for applications in motivic tensor-triangular geometry. These are supplied contextual claims, not independently verified results in this draft.
\subsection{Short English statement}
The proposed conjecture says that the periodic points in the tensor-triangular spectrum of compact derived permutation modules are exactly the nonclosed points. The supplied passage reports the finite \(p\)-group and pro-\(p\) cases as established and labels the general problem partially solved. A precise standalone statement still requires the coefficient hypotheses, the definition of periodicity, and the meaning and parameter range of \(P(H,p)\); the source's two uses of \(p\) cannot safely be identified.
\subsection{Sources}
\begin{itemize}
\item Supplied local UnsolvedMath registry record, permanent ID \(30006651\), source problem code \texttt{catalogue entry}, title \emph{Periodic Loci for Derived Permutation Modules}, in the Oberwolfach Reports Open Problems collection. The record supplies Conjecture 7, Remark 8, and Theorem 9 as the relevant passage. No individual proposer is supplied.
\item \emph{Cohomology of Finite Groups: Interactions and Applications} (2026), Oberwolfach Reports, pp.\ 541--610, DOI \texttt{10.4171/owr/2026/9}. Supplied originating report link: \url{https://doi.org/10.4171/owr/2026/9}.
\item Miller, \emph{Permutation twisted cohomology, remixed}. Supplied literature link: \url{https://arxiv.org/abs/2509.00954}. The excerpt uses the reference label [Mil25] for Miller's work.
\item Balmer--Gallauer, \emph{The tt-geometry of permutation modules}. Supplied literature link: \url{https://arxiv.org/abs/2307.04398}. The excerpt separately uses [BG25]; the supplied record does not give its full bibliography or explicitly establish the correspondence between that label and this linked reference.
\item The excerpt attributes Theorem 9 and motivic applications to [Gal25]. No full bibliographic entry or URL for [Gal25] is supplied.
\item Literature assessment date: \texttt{2026-08-22}. This is supplied dated curation, not independently verified. It labels the problem partially solved, reports the pro-\(p\) and important elementary abelian cases as established, and describes the general profinite case as open because the needed ampleness is not known. The background expressly requests independent verification of time-sensitive status and cited claims.
\item The registry labels the statement \texttt{corrected\_verified} and says it was reconstructed and checked against the official report. This is supplied provenance metadata; this draft does not independently verify that check or claim to have visited any linked source.
\end{itemize}
\end{document}
